\paragraph{厚生についての順位}
\label{sec:chapter5_beta_shock_welfare}
次に価格分散\(\Delta_t^H\)を考慮した場合を例として
どの政策が最も厚生\(W_s^H\)の下落を抑制できているかを確認する。
厚生\(W_s^H\)の大きい順に政策を示す。
\begin{table}[H]
\centering
\caption{厚生についての順位}
\label{tab:chapter5_beta_shock_welfare}
\begin{tabular}{lc}
\hline
金融政策 & \(W_s^H\)（\(\Delta_t^H\)有り） \\
\hline
名目総消費水準目標（NCLT） & $-1.1677$ \\
潜在生産水準目標（POLT） & $-1.2049$ \\
生産水準目標（OLT） & $-1.2141$ \\
実質消費水準目標（CLT） & $-1.2332$ \\
名目GDP水準目標（NGDPLT） & $-1.2797$ \\
消費者物価テイラールール（TR-CPI） & $-1.3768$ \\
生産者物価水準目標（PPLT） & $-1.4120$ \\
生産者物価テイラールール（TR-PPI） & $-1.6796$ \\
生産者物価インフレ目標（IT-PPI） & $-1.8149$ \\
消費者物価水準目標（CPLT） & $-2.4896$ \\
消費者物価インフレ目標（IT-CPI） & $-2.4957$ \\
\hline
\end{tabular}
\end{table}
NCLTが最も高い厚生\(W_s^H\)を達成しており\(\beta\)ショックによる負の影響を最小限に留めていることがわかる。
一方で消費者物価指数（CPI）\(p_t^{H\to W}\)のみを目標に含むIT-CPIやCPLTは極めて低い評価となった。